\documentclass[12pt,a4paper]{report}
\usepackage[french]{babel}
\usepackage[T1]{fontenc}
\usepackage{times}
\usepackage{geometry}
\usepackage{xcolor}
\usepackage{graphicx}
\usepackage{hyperref}
\usepackage{amsmath}
\usepackage{listings}
\usepackage{longtable}
\usepackage{float}

% ===== AJOUT POUR PIED DE PAGE =====
\usepackage{fancyhdr}
\usepackage{lastpage}

\geometry{left=2.5cm, right=2.5cm, top=2.5cm, bottom=2.5cm}

% Style des titres (sobres, sans couleur)
\usepackage{titlesec}
\titleformat{\chapter}{\normalfont\fontsize{16}{20}\bfseries}{\thechapter}{1em}{}
\titleformat{\section}{\normalfont\fontsize{14}{18}\bfseries}{\thesection}{1em}{}
\titleformat{\subsection}{\normalfont\fontsize{13}{16}\bfseries}{\thesubsection}{1em}{}

% ===== CONFIGURATION DU PIED DE PAGE =====
\pagestyle{fancy}
\fancyhf{}  % Efface tous les en-têtes et pieds existants

% Pied de page
\fancyfoot[L]{\small \textit{LocaleGuide - Projet WEB}}  % Gauche
\fancyfoot[C]{\thepage}  % Centre (numéro de page)
 % Droite (date)

% Épaisseur du trait au-dessus du pied de page
\renewcommand{\footrulewidth}{0.4pt}

% Style pour les premières pages de chapitre (sans en-tête mais avec pied)
\fancypagestyle{plain}{
    \fancyhf{}
    \fancyfoot[C]{\thepage}
    \renewcommand{\headrulewidth}{0pt}
    \renewcommand{\footrulewidth}{0.4pt}
}

\usepackage{longtable}
\usepackage{booktabs}
\usepackage{array}
\usepackage{listings}
\usepackage{xcolor}
\usepackage{listings}
\usepackage{xcolor}

\lstset{
  basicstyle=\ttfamily\small,
  keywordstyle=\color{black}\bfseries,
  stringstyle=\color{black},
  commentstyle=\color{gray}\itshape,
  numbers=left,
  numberstyle=\tiny\color{gray},
  stepnumber=1,
  numbersep=5pt,
  breaklines=true,
  frame=single,
  rulecolor=\color{black!30},
  backgroundcolor=\color{gray!5},
  tabsize=2,
  showstringspaces=false,
  language=SQL
}


\begin{document}

% ========== PAGE DE GARDE ==========
\begin{titlepage}
    \noindent
    \begin{minipage}{0.3\textwidth}
        \includegraphics[width=2.5cm]{isimmlogo.jpg}
    \end{minipage}
    \hfill
    \begin{minipage}{0.3\textwidth}
        \raggedleft
        \includegraphics[width=2.5cm]{iso.png}
    \end{minipage}
    
    \vspace{1cm}
    
    \centering
    {\color{black}\fontsize{20}{24}\bfseries LocaleGuide\par}
    \vspace{0.5cm}
    
    {\color{gray}\fontsize{14}{18}\bfseries Plateforme de Local Guide \\ pour Monastir Centre\par}
    \vspace{2.5cm}
    
    {\color{darkgray}\fontsize{14}{16}\textbf{Présenté par :}\par}
    \vspace{0.3cm}
    {\color{black}\fontsize{12}{14} Mejri Maram\par}
    {\color{black}\fontsize{12}{14} Werghemmi Isra\par}
    {\color{black}\fontsize{12}{14} Mejbri Tasnim\par}
    {\color{black}\fontsize{12}{14} Louati Souha\par}
    \vspace{2cm}
    
    {\color{darkgray}\fontsize{14}{16}\textbf{Enseignante :}\par}
    \vspace{0.3cm}
    {\color{black}\fontsize{12}{14}Senda Laatawi\par}
    \vspace{1.5cm}
    
    {\color{darkgray}\fontsize{14}{16}\textbf{Établissement :}\par}
    \vspace{0.3cm}
    {\color{black}\fontsize{12}{14}Institut Supérieur d'Informatique et de Mathématiques\\ Monastir\par}
    \vspace{1.5cm}
    
    {\color{darkgray}\fontsize{14}{16}\textbf{Année universitaire :}\par}
    \vspace{0.3cm}
    {\color{black}\fontsize{12}{14}2026--2026\par}
    
    \vfill
\end{titlepage}

\tableofcontents
\newpage

% ========== CHAPITRE 1 : INTRODUCTION ==========
\chapter{Introduction et présentation de l’idée}

\section{Contexte et objectifs}

Dans un contexte de digitalisation du secteur touristique, les utilisateurs recherchent des solutions numériques fiables leur permettant de découvrir facilement des commerces locaux, des guides touristiques ainsi que des événements.

Le projet \textbf{LocalGuide Monastir} s’inscrit dans cette perspective en proposant une application web moderne permettant de centraliser ces services sur une seule plateforme.

Les objectifs principaux du projet sont :
\begin{itemize}
    \item Faciliter la découverte des commerces locaux
    \item Mettre en relation les utilisateurs avec des guides touristiques
    \item Permettre la consultation et la publication d’avis
    \item Offrir un système de réservation simple et efficace
\end{itemize}


\section{L’idée originale}

L’idée principale du projet consiste à développer une plateforme interactive qui connecte différents acteurs du domaine touristique : les utilisateurs, les commerces locaux et les guides.

L’application permet :
\begin{itemize}
    \item La recherche intelligente de commerces et services
    \item L’ajout de nouveaux commerces par les utilisateurs
    \item La candidature pour devenir guide touristique
    \item La gestion dynamique des contenus via une base de données
\end{itemize}

Cette solution apporte une valeur ajoutée en centralisant les informations et en améliorant l’expérience utilisateur.

\begin{itemize}
  \item \textbf{Découvrir} les commerces locaux (restaurants, cafés, boutiques, artisans)
        avec leurs horaires, prix et avis clients.
  \item \textbf{Réserver} des guides touristiques certifiés, spécialisés dans différents
        thèmes (gastronomie, histoire, mer, artisanat, nocturne).
  \item \textbf{Suivre} les actualités locales : événements, promotions, nouvelles ouvertures.
  \item \textbf{Contribuer} à la communauté en laissant des avis ou en proposant ses
        services de guide.
\end{itemize}

\section{Défense de l’idée (justification)}

Le choix de ce projet repose sur plusieurs raisons :

\begin{itemize}
    \item \textbf{Utilité réelle} : répondre aux besoins des touristes et des habitants
    \item \textbf{Innovation} : intégration de fonctionnalités modernes (recherche avancée, avis, réservation)
    \item \textbf{Accessibilité} : interface simple et intuitive
    \item \textbf{Impact local} : valorisation des commerces et guides locaux
\end{itemize}

Ce projet constitue ainsi une solution pertinente, à la fois utile et évolutive, dans le domaine du tourisme numérique.

% ========== CHAPITRE 2 : CONCEPTION ==========
\chapter{Conception}
\section{Diagramme des cas d'utilisation}
\subsection*{ Diagramme des cas d'utilisation principal}

Le système \textbf{LocalGuide} implique quatre acteurs principaux :

\begin{itemize}
    \item \textbf{Visiteur} (non connecté) : peut consulter les commerces, les guides et les actualités.
    \item \textbf{Utilisateur} (connecté) : peut laisser des avis, réserver des guides et des commerces.
    \item \textbf{Candidat Guide} : peut soumettre une demande pour devenir guide.
    \item \textbf{Administrateur} : modère les avis, valide les guides et les commerces, et gère toutes les actualités.
\end{itemize}
\begin{figure}[h]
    \centering
    \includegraphics[width=0.8\textwidth]{usecase.png}
    \caption{le diagramme des cas d'utilisation du système LocalGuide.}
    \label{fig:monimage}
\end{figure}

La figure~\ref{fig:usecase} présente le diagramme des cas d'utilisation du système LocalGuide.

\section {Détail des cas d'utilisation}
\subsection {Détail des cas d’utilisation (UC détaillés)}

\subsubsection{UC1 : Ajouter un commerce}

\textbf{Acteur principal :} Utilisateur

\textbf{Préconditions :}
\begin{itemize}
    \item L’utilisateur est authentifié
\end{itemize}

\textbf{Postconditions :}
\begin{itemize}
    \item Le commerce est enregistré dans la base avec statut \textit{pending} (en attente de validation admin)
\end{itemize}

\textbf{Scénario nominal :}
\begin{enumerate}
    \item L’utilisateur accède à la page ``Ajouter un commerce''
    \item Il remplit les informations (nom, catégorie, description, adresse, image, etc.)
    \item Il soumet le formulaire
    \item Le système envoie les données au backend (Supabase)
    \item Le commerce est enregistré dans la table \texttt{commerces}
    \item Un message de confirmation s’affiche
\end{enumerate}

\textbf{Scénarios alternatifs :}
\begin{itemize}
    \item Si un champ obligatoire est vide $\rightarrow$ message d’erreur
    \item Si erreur serveur $\rightarrow$ message ``échec d’ajout''
\end{itemize}

\vspace{0.5cm}

\subsubsection{UC2 : Devenir guide}

\textbf{Acteur principal :} Utilisateur

\textbf{Préconditions :}
\begin{itemize}
    \item L’utilisateur est authentifié
\end{itemize}

\textbf{Postconditions :}
\begin{itemize}
    \item Une demande de guide est enregistrée avec statut \textit{pending}
\end{itemize}

\textbf{Scénario nominal :}
\begin{enumerate}
    \item L’utilisateur accède à la page ``Devenir guide''
    \item Il remplit le formulaire (nom, spécialité, langues, expérience, etc.)
    \item Il télécharge sa photo
    \item Il valide la demande
    \item Le système envoie les données au backend
    \item Les données sont insérées dans la table \texttt{guides}
    \item Un message de confirmation s’affiche
\end{enumerate}

\textbf{Scénarios alternatifs :}
\begin{itemize}
    \item Si les données sont invalides $\rightarrow$ message d’erreur
    \item Si l’upload de l’image échoue $\rightarrow$ demande de réessayer
\end{itemize}


\section{Diagrammes de séquence}

\subsection{Séquence : Réserver un guide}

\begin{figure}[H]
    \centering
    \includegraphics[width=0.8\textwidth]{guide_sequence.png}
    \caption{Diagramme de séquence - Réservation d'un guide}
    \label{fig:sequence_reservation}
\end{figure}

\subsection{Séquence : Authentification et création de profil}

\begin{figure}[H]
    \centering
    \includegraphics[width=0.8\textwidth]{sequence.png}
    \caption{Diagramme de séquence - Authentification et création de profil}
    \label{fig:sequence_auth}
\end{figure}

\section{Diagramme entité-association (modèle de données)}

Le schéma relationnel de LocalGuide s'articule autour de six tables principales :
\begin{figure}[h]
    \centering
    \includegraphics[width=0.8\textwidth]{diagramme_classe.png}
    \caption{diagramme de classe}
    \label{fig:monimage}
\end{figure}
\section{Architecture technique}

L’architecture du système repose sur une architecture en trois couches :

\begin{itemize}
    \item \textbf{Frontend :} HTML / CSS / JavaScript
    \item \textbf{Backend :} Supabase (API + Auth + logique)
    \item \textbf{Base de données :} PostgreSQL (via Supabase)
\end{itemize}

\textbf{Flux de données :}

\begin{center}
Frontend $\rightarrow$ API Supabase $\rightarrow$ Base de données
\end{center}

\section{Maquettes / wireframes}
\begin{figure}[h]
    \centering
    \includegraphics[width=0.8\textwidth]{accueil.png}
    \caption{accueil}
    \label{fig:monimage}
\end{figure}
\begin{figure}[h]
    \centering
    \includegraphics[width=0.8\textwidth]{commerce.png}
    \caption{commerce }
    \label{fig:monimage}
\end{figure}
\begin{figure}[h]
    \centering
    \includegraphics[width=0.8\textwidth]{nouveauté.png}
    \caption{nouveauté}
    \label{fig:monimage}
\end{figure}
\begin{figure}[h]
    \centering
    \includegraphics[width=0.8\textwidth]{guides.png}
    \caption{guides}
    \label{fig:monimage}
\end{figure}
\begin{figure}[h]
    \centering
    \includegraphics[width=0.8\textwidth]{admin.png}
    \caption{Admin}
    \label{fig:monimage}
\end{figure}
\begin{figure}[h]
    \centering
    \includegraphics[width=0.8\textwidth]{ajout commerce.png}
    \caption{ajout commerce}
    \label{fig:monimage}
\end{figure}

% ========== CHAPITRE 3 : DÉVELOPPEMENT ==========
\chapter{Développement}

\section{Technologies choisies et justifications}
% Votre texte ici
\begin{itemize}
  \item \textbf{Front-end} : HTML5 / CSS3 / JS ES6
    \begin{itemize}
      \item Légèreté, pas de build tool nécessaire, déploiement immédiat
    \end{itemize}
  
  \item \textbf{BaaS} : Supabase
    \begin{itemize}
      \item PostgreSQL managé + Auth + RLS + Realtime + SDK JS intégré
    \end{itemize}
  
  \item \textbf{Base de données} : PostgreSQL 15
    \begin{itemize}
      \item Robustesse, extensions PostGIS et pg\_trgm, fonctions RPC
    \end{itemize}
  
  \item \textbf{Recherche} : pg\_trgm + tsvector
    \begin{itemize}
      \item Recherche full-text en français avec tolérance aux fautes
    \end{itemize}
  
  \item \textbf{Géolocalisation} : PostGIS
    \begin{itemize}
      \item Requêtes de proximité spatiale (\texttt{ST\_DWithin})
    \end{itemize}
  
  \item \textbf{Sécurité} : Row Level Security
    \begin{itemize}
      \item Contrôle d'accès granulaire directement dans la base de données
    \end{itemize}
\end{itemize}

\section{Organisation du code}


\begin{lstlisting}[basicstyle=\ttfamily\fontsize{12}{14}\selectfont, numbers=none, frame=none, xleftmargin=0pt]
LocalGuide/
│
├── .github/workflows/ci.yml
│
├── backend/supabase/
│   ├── migrations/          (Scripts de migration DB)
│   ├── seeds/               (Donnees initiales)
│   └── functions/           (Fonctions PostgreSQL)
│
├── docs/
│   ├── rapport/             (Rapport du projet)
│   └── specifications/      (Cahier des charges)
│
├── frontend/
│   ├── css/                 (Styles CSS)
│   ├── images/              (Images et logos)
│   └── assets/              (Autres ressources)
│
├── js/
│   ├── supabaseClient.js    (Connexion Supabase)
│   ├── common.js            (Fonctions partagees)
│   ├── accueil.js            (Page d'accueil)
│   ├── fichecomm.js         (Gestion commerces)
│   ├── guide.js             (Reservation guides)
│   ├── news.js              (Actualites)
│   ├── admin.js             (Dashboard admin)
│   └── add_news.js          (Ajout actualites)
│
├── pages/
│   ├── accueil.html          (Accueil)
│   ├── fichecomm.html       (Annuaire commerces)
│   ├── guide.html           (Liste guides)
│   ├── news.html            (Actualites)
│   ├── add.html             (Ajout commerce)
│   ├── add_news.html        (Ajout actualite)
│   ├── avis.html            (Deposer avis)
│   ├── devenir_guide.html   (Candidature)
│   ├── login.html           (Connexion)
│   ├── Register.html        (Inscription)
│   └── admin.html           (Administration)
│
├── public/
│   ├── css/                 (Styles publics)
│   └── images/              (Images publiques)
│
├── scripts/
│   ├── deploy.sh            (Script deploiement)
│   └── backup.sh            (Script sauvegarde)
│
├── src/                     (Sources additionnelles)
│
├── supabase/
│   ├── config.toml          (Configuration)
│   └── seed.sql             (Donnees de test)
│
├── tests/
│   ├── unit/                (Tests unitaires)
│   └── integration/         (Tests integration)
│
├── .env.example             (Variables d'environnement)
├── .gitignore               (Fichiers ignores Git)
├── arborescence.txt         (Arborescence generee)
├── index.html               (Page principale)
├── package-lock.json        (Verrouillage versions)
└── package.json             (Dependances Node.js)
\end{lstlisting}

\section{Implémentation des fonctionnalités clés}
\subsection{Authentification et gestion des rôles}

La connexion utilise le SDK Supabase Auth. Après authentification, le rôle de l'utilisateur
est lu depuis la table \texttt{users\_profiles} pour adapter la navigation (affichage
conditionnel du lien Admin).
\begin{lstlisting}[caption=Connexion et détection du rôle (login.html)]

const { data, error } = await supabase.auth.signInWithPassword({
            email: email,
            password: password
        })
        
        if (error) {
            messageDiv.innerHTML = 'erreur ' + error.message
            messageDiv.className = 'message error'
            return
        }
        
        // Déterminer le rôle
        let role = 'client'
        const userEmail = data.user.email
        
        //  FORCER ADMIN POUR CET EMAIL
        if (userEmail === 'admin@localguide.com') {
            role = 'admin'
            console.log(' Admin détecté !')
        } else {
            // Vérifier dans users_profiles
            const { data: profile } = await supabase
                .from('users_profiles')
                .select('role')
                .eq('id', data.user.id)
                .single()
            
            if (profile && profile.role) {
                role = profile.role
            }
        }
        
        messageDiv.innerHTML = ` Connecté en tant que ${role} ! Redirection...`
        messageDiv.className = 'message success'
        
        // Rediriger
        setTimeout(() => {
            if (role === 'admin') {
                window.location.href = 'admin.html'
            } else {
                window.location.href = 'accueil.html'
            }
        }, 1000)
\end{lstlisting}

\subsection{Recherche de commerces}
% Votre code ici
\begin{lstlisting}[caption=Rechercher des commerces (fichecomm.html)]

let liste = [..allCommerces]
    
    // Filtre par catégorie
    if (activeFilter !== "all") {
        liste = liste.filter(c => c.categorie === activeFilter)
    }
    
    //  FILTRE DE RECHERCHE TEXTE
    if (searchQuery.trim()) {
        const q = searchQuery.toLowerCase()
        liste = liste.filter(c =>
            c.nom.toLowerCase().includes(q) ||
            c.description?.toLowerCase().includes(q)
        )
    }
    // Tri
    if (sortBy === "note") liste.sort((a, b) => b.note - a.note)
    if (sortBy === "nom") liste.sort((a, b) => a.nom.localeCompare(b.nom))
    
    displayedCommerces = liste
    currentPage = 1
    hasMore = displayedCommerces.length > itemsPerPage
    
    const countDisplay = document.getElementById("countDisplay")
    if (countDisplay) countDisplay.textContent = displayedCommerces.length
    
    renderCurrentPage()
    renderPaginationButtons()
}
\end{lstlisting}

\subsection{Gestion des guides}
\begin{lstlisting}[ caption=exemlpe de code de gestion guides(guide.js)]
async function loadGuidesEnAttente() {
    const { data, error } = await supabase
        .from('guides')
        .select('*')
        .eq('statut', 'inactif')
        .order('created_at', { ascending: false })
    
    if (error) {
        console.error('Erreur chargement guides:', error)
        return []
    }
    return data
}

async function approveGuide(id) {
    if (!confirm(' Approuver ce guide ? Il sera visible sur la page des guides.')) return
    
    const { error } = await supabase
        .from('guides')
        .update({ 
            statut: 'disponible',
            updated_at: new Date()
        })
        .eq('id', id)
    
    if (error) {
        alert(' Erreur: ' + error.message)
    } else {
        alert('Guide approuvé ! Il est maintenant visible sur le site.')
        await renderGuidesEnAttente()
    }
}


\end{lstlisting}
\subsection{Upload de la photo du guide}
\begin{lstlisting}
async function handleFile(file) {
    const error = document.getElementById("pictureError");
    const maxSize = 5 * 1024 * 1024;
    const allowed = ["image/jpeg", "image/png", "image/webp"];

    if (!allowed.includes(file.type)) {
      showError(error, "Format non supporté. Utilisez JPG, PNG ou WEBP.");
      return;
    }
    
    if (file.size > maxSize) {
      showError(error, `Fichier trop lourd (max 5MB).`);
      return;
    }

    try {
      const fileExt = file.name.split('.').pop();
      const fileName = `guide_${Date.now()}_${Math.random().toString(36).substring(7)}.${fileExt}`;
      
      // Upload vers bucket 'guides-photos'
      const { error: uploadError } = await supabase.storage
        .from('guides-photos')
        .upload(fileName, file)
      
      if (uploadError) throw uploadError
      
      // Récupérer l'URL publique
      const { data: { publicUrl } } = supabase.storage
        .from('guides-photos')
        .getPublicUrl(fileName)
      
      uploadedImageUrl = publicUrl
      
    } catch (err) {
      console.error('Erreur:', err);
      showError(error, "Erreur inattendue lors de l'upload");
    }
}
\end{lstlisting}
\subsection{Sécurité Row Level Security}
Toutes les tables sont protégées par des politiques RLS strictes :

\begin{lstlisting}[language=SQL, captionpos=b, caption=Exemple de politiques RLS sur la table guides]
-- Tout le monde peut lire les guides
CREATE POLICY "public_view_guides"
  ON public.guides FOR SELECT TO public USING (true);

-- Seul le proprietaire ou un admin peut modifier
CREATE POLICY "owner_or_admin_update_guides"
  ON public.guides FOR UPDATE TO authenticated
  USING (auth.uid() = user_id OR auth.uid() IN (
      SELECT user_id FROM public.user_roles WHERE role = 'admin'
  ));
\end{lstlisting}


% ========== CHAPITRE 4 : TESTS ==========
\chapter{Tests et validation}

\section{Environnement et outils de test}

Les tests ont été effectués dans l'environnement suivant :

\begin{itemize}
  \item \textbf{Navigateurs testés} : Google Chrome (v124), Brave et le navigateur intégré d'Android.
  \item \textbf{Base de données} : Console Supabase (éditeur SQL intégré) pour la vérification des requêtes et des politiques RLS (\textit{Row Level Security}).
  \item \textbf{Tests d'interface} : Tests manuels des formulaires de connexion, d'inscription, d'ajout de commerce et de réservation.
  \item \textbf{Validation des données} : Vérification des contraintes \texttt{NOT NULL}, des types et des clés étrangères directement via Supabase.
  \item \textbf{Tests géographiques} : Vérification de la localisation des commerces et des guides sur carte (coordonnées \texttt{latitude}/\texttt{longitude}).
\end{itemize}

\section{Scénarios de test}

\begin{table}[H]
\centering

\begin{tabular}{|p{5.5cm}|c|p{5.5cm}|}
\hline
\textbf{Fonctionnalité} & \textbf{Statut} & \textbf{Observations} \\
\hline
Inscription utilisateur & OK & Création du profil dans \texttt{users\_profiles} \\
\hline
Connexion utilisateur & OK & Authentification via Supabase Auth \\
\hline
Navigation sans connexion (visiteur) & OK & Accès en lecture seule aux commerces \\
\hline
Ajout d'un commerce (guide) & OK & Statut initial \texttt{en\_attente} \\
\hline
Validation d'un commerce (admin) & OK & Passage du statut à \texttt{approuvé} \\
\hline
Refus d'un commerce avec motif (admin) & OK & Champ \texttt{raison\_refus} renseigné \\
\hline
Recherche de commerces par nom & OK & Recherche insensible à la casse \\
\hline
Affichage des commerces sur carte & OK & Coordonnées GPS correctement récupérées \\
\hline
Dépôt d'un avis sur un commerce & OK & Note globale et notes détaillées enregistrées \\
\hline
Réservation auprès d'un guide & OK & Statut initial \texttt{en\_attente} \\
\hline
Confirmation de réservation (guide) & OK & Statut mis à jour en \texttt{confirmée} \\
\hline
Annulation de réservation & OK & Statut mis à jour en \texttt{annulée} \\
\hline
Publication d'une actualité & OK & Statut \texttt{publié} après validation \\
\hline
Affichage du profil guide & OK & Informations et spécialités affichées \\
\hline
Gestion des rôles (admin) & OK & Attribution correcte via \texttt{user\_roles} \\
\hline
\end{tabular}

\end{table}


% ========== CHAPITRE 5 : CONCLUSION ==========
\chapter{Conclusion et perspectives}

\section{Bilan du projet}

Le projet a permis de développer une application complète intégrant frontend et backend réel, de manipuler une base de données dynamique.

\section{Difficultés rencontrées et apprentissages}

\begin{itemize}
  \item \textbf{Row Level Security} : La mise en place des politiques RLS a demandé
        une compréhension fine des rôles Supabase et de la différence entre
        \texttt{TO public} et \texttt{TO authenticated}.
  \item \textbf{Trigger de création de profil} : La synchronisation entre
        \texttt{auth.users} et \texttt{users\_profiles} via trigger PostgreSQL a
        nécessité plusieurs itérations pour gérer les cas d'erreur.
  \item \textbf{PostGIS} : L'activation de l'extension et la compréhension des
        types \texttt{GEOMETRY(Point, 4326)} ont représenté une courbe d'apprentissage
        non négligeable.
  \item \textbf{Modules ES6} : La gestion des imports JavaScript entre fichiers
        partagés (\texttt{common.js}, \texttt{supabaseClient.js}) en mode module
        a requis une organisation rigoureuse.
\end{itemize}

\section{Évolutions possibles}

\begin{itemize}
  \item \textbf{Application mobile} : portage en Progressive Web App (PWA) ou
        application React Native.
  \item \textbf{Système de notifications} : push notifications via Supabase Realtime
        pour les nouvelles réservations et validations.
  \item \textbf{Paiement en ligne} : intégration d'une passerelle de paiement
        (Stripe ou solution locale tunisienne) pour les réservations guides.
  \item \textbf{Carte interactive} : affichage cartographique des commerces et guides
        via Leaflet.js en exploitant les données PostGIS.
  \item \textbf{Extension géographique} : déploiement de la plateforme à d'autres
        villes tunisiennes (Sousse, Sfax, Tunis).
  \item \textbf{Tableau de bord analytique} : statistiques de fréquentation,
        notes moyennes, tendances de réservation pour les commerçants et guides.
\end{itemize}

% ========== ANNEXES ==========
\chapter{Annexes}

\section{Lien vers le dépôt Git}
\url{https://github.com/Souha-05/local}

\section*{Captures d'écran}

\begin{figure}[H]
    \centering
    \includegraphics[width=0.9\textwidth]{accueil.png}
    \caption{Page d'accueil LocalGuide — hero Monastir}
    \label{fig:accueil}
\end{figure}

\begin{figure}[H]
    \centering
    \includegraphics[width=0.9\textwidth]{admin.png}
    \caption{Tableau de bord administrateur — modération commerces}
    \label{fig:admin}
\end{figure}

\begin{figure}[H]
    \centering
    \includegraphics[width=0.9\textwidth]{guide.png}
    \caption{Page guides — cartes avec filtres par spécialité}
    \label{fig:guides}
\end{figure}

\subsection{Fichiers de configuration}

\subsubsection{Configuration du client Supabase}

Le fichier de configuration principal pour la connexion à Supabase est présenté ci-dessous. Il utilise les variables d'environnement pour sécuriser les identifiants.

\begin{lstlisting}[language=JavaScript,  captionpos=b,caption={supabase.js - Connexion à Supabase}]
const supabaseUrl = import.meta.env.VITE_SUPABASE_URL
const supabaseAnonKey = import.meta.env.VITE_SUPABASE_ANON_KEY

import { createClient } from 'https://cdn.jsdelivr.net/npm/@supabase/supabase-js@2/+esm'

export const supabase = createClient(supabaseUrl, supabaseAnonKey)

console.log(' Supabase connecté')
\end{lstlisting}

\subsubsection{Fichier package.json}

Le fichier \texttt{package.json} définit les dépendances et les scripts nécessaires au lancement de l'application front-end.

\begin{lstlisting}[language=JSON, captionpos=b,caption={package.json - Configuration Vite}]
{
  "name": "web",
  "version": "0.0.0",
  "private": true,
  "type": "module",
  "scripts": {
    "dev": "vite",
    "build": "vite build",
    "preview": "vite preview"
  },
  "devDependencies": {
    "vite": "^8.0.8"
  }
}
\end{lstlisting}

\subsection{Fichier d'environnement (.env)}

Les identifiants de connexion à Supabase sont stockés dans un fichier \texttt{.env} à la racine du projet (non versionné pour des raisons de sécurité) :

\begin{lstlisting}[language=bash, captionpos=b, caption={.env - Variables d'environnement}]
VITE_SUPABASE_URL=https://auumbkcfwjpvdifnxzly.supabase.co
VITE_SUPABASE_ANON_KEY=sb_publishable_WXvPDV94YwVHQGAuVK0Nqg_Z3rIk1it
\end{lstlisting}

\end{document}